\section{Tests et validation}
\label{sec:tests}

% Évaluation selon les caractéristiques de qualité de la norme ISO/IEC 25010.

\subsection{Stratégie de test}
% Décrivez votre approche : tests unitaires, d'intégration, fonctionnels,
% de performance, de non-régression, tests manuels...
% Précisez les outils utilisés (pytest, JUnit, Selenium...).

\subsection{Cas de tests}
% Listez les cas de test principaux.
% Format : identifiant, description, résultat attendu, statut.

\begin{center}
  \begin{tabular}{llll}
    \hline
    \textbf{ID} & \textbf{Description} & \textbf{Résultat attendu} & \textbf{Statut} \\
    \hline
    T-01 & ... & ... & Succès \\
    T-02 & ... & ... & Succès \\
    \hline
  \end{tabular}
\end{center}

\subsection{Couverture qualité (ISO/IEC 25010)}
% Évaluez les caractéristiques de qualité pertinentes pour votre projet.
% Traitez les critères applicables parmi : fonctionnalité, fiabilité,
% performance, maintenabilité, sécurité, portabilité.

\subsection{Résultats}
% Taux de réussite global, bugs identifiés et corrigés,
% problèmes en suspens et impact sur le projet.
